\documentclass[10pt,a4paper]{amsart}
\usepackage[margin=19mm]{geometry}
\usepackage{amsmath,amssymb,mathtools}
\usepackage[scr=rsfs,cal=euler]{mathalfa}
\usepackage{xfrac}
\usepackage[unicode,hidelinks,bookmarks=false]{hyperref}
\newcommand{\ie}{i.e.\ }
\newcommand{\cf}{cf.\ }
\newcommand{\resp}{respectively\ }
\newcommand{\Subsection}{\subsection}
\providecommand{\nobreakdash}{\nobreak\hbox{-}}
\setlength{\emergencystretch}{2em}
\multlinegap0pt
\usepackage{amssymb}
\newtheorem{lemma}{Lemma}
\newtheorem{proposition}{Proposition}
\usepackage{cleveref}
\crefname{lemma}{Lemma}{Lemmas}
\crefname{proposition}{Proposition}{Propositions}
\newcommand{\mM}{\mathcal{M}}
\title{On various notions of distance between subalgebras of operator algebras}
\author{Ved Prakash Gupta, Sumit Kumar}
\date{Source: arXiv:2403.05967v1 (CC BY 4.0)}
\begin{document}
\maketitle

\noindent\emph{Source and licence.} On various notions of distance between subalgebras of operator algebras. Version arXiv:2403.05967v1 (CC BY 4.0), distributed by arXiv under the Creative Commons Attribution 4.0 International licence, http://creativecommons.org/licenses/by/4.0/, as recorded in the arXiv licence metadata for this exact version and shown on the abstract page https://arxiv.org/abs/2403.05967v1. Full licence notice: \texttt{licenses/arxiv-2403.05967-CC-BY-4.0.txt}.\par\smallskip

\noindent\textit{Editorial scope.} Counted unit: the article's proposition dMT-SOT (source0.tex lines 1496--1504). The article's complete original proof of it is delivered with it (source0.tex lines 1505--1579, one \texttt{proof} environment of 75 lines). No mathematics is written, repaired or re-derived: the statement and the proof are the article's own lines, in the article's own notation, and no retained line is substituted or re-flowed.  Retained uncounted as the source-local context the proof needs: lines 1458--1470.  Nothing was omitted from any retained span, so the package carries no omission record.  No retained line contains a citation, so the wrapper carries no bibliography and no bibliography entry.  No retained line contains a figure environment, an image-inclusion command or drawing code, so no image is delivered and the package declares no asset dependency.  Editorial formatting only, in addition: an AMS \texttt{amsart} preamble that restates the article's own declarations and macros used by the retained lines, namely the environments \texttt{lemma}, \texttt{proposition} on separate counters, together with the standard packages those lines need.  The retained environments print in this document as Lemma 1, Proposition 1 in order of appearance, whereas the article prints the same items with the numbering of its own full document.  The complete native source of the article as distributed in the arXiv e-print for this exact version is bundled unchanged as \texttt{sources/155839/source0.tex}.
\begin{lemma}\label{dMT-facts}
  Let $\mM$ be a von Neumann algebra with a  faithful normal
  tracial state $\tau$. Then, the following hold:
\begin{enumerate}
\item  $d_{MT}$  is a semi-metric on $\mathrm{Sub}_{\mM}$.

\item $d_{MT}$ is a metric on $W^*\text{-}\mathrm{Sub}_{\mM}$. 

\item \(
d_C(P,Q)\leq d_{MT}(P,Q) 
\)    for all $P, Q \in \text{Sub}_{\mM}$.
 \end{enumerate}
\end{lemma}

\begin{proposition}\label{dMT-SOT}
Let $(\mM, \tau)$ be as in \Cref{dMT-facts}. Then, 
\[
d_{MT}(P,Q) = d_{MT}\Big(P, \overline{Q}^{S.O.T.}\Big) =
d_{MT}\Big(\overline{P}^{S.O.T.}, \overline{Q}^{S.O.T.}\Big) = d_{MT}\Big(P'',
Q''\Big)
\] 
for all $P, Q \in *\text{-}\mathrm{Sub}^u_{\mM}$.
\end{proposition}

\begin{proof}
In order to avoid any possible confusion, for every $x \in \mM$ and $S
\subseteq \mM$, let $d_\tau\big(\widehat{x}, \widehat{S}\big):=\inf \{ \|\widehat{x} -
\widehat{s}\|_\tau : s \in S\}$.

Now, let $P, Q \in *\text{-Sub}^u_{\mM}$.  We first assert that
$d_\tau\Big(\widehat{a}, \widehat{B_1(Q)}\Big) = d_\tau\Big(\widehat{a},
\widehat{B_1(\overline{Q}^{S.O.T.})}\Big)$ for all $a \in B_1(P)$.

Clearly, $d_\tau\Big(\widehat{a}, \widehat{ B_1(\overline{Q}^{S.O.T.}\Big)}\leq
d_\tau\Big(\widehat{a}, \widehat{B_1(Q)}\Big)$ for all $a\in B_1(P)$. For the
reverse inequality, fix an $a \in B_1(P)$ and an $\epsilon>0$. Then,
there exists a $b_{0}\in B_1(\overline{Q}^{S.O.T.})$ such that
\[
d_\tau\Big(\widehat{a},
\widehat{B_1(\overline{Q}^{S.O.T.})}\Big) \leq
\|\widehat{a}-\widehat{b_{0}}\|_{\tau}< d_\tau\Big(\widehat{a},
\widehat{B_1(\overline{Q}^{S.O.T.})}\Big)+ \epsilon.
\]
As $ B_1\Big(\overline{Q}^{S.O.T.}\Big)= \overline{B_1(Q)}^{S.O.T.}$ (by
Kaplansky density theorem), there exists a net
$\{b_{\alpha}\}\subseteq B_1(Q)$ such that
$b_{\alpha}\stackrel{S.O.T.}{\longrightarrow} b_{0}$. Hence, there
exists an $\alpha_{0}$ such that
\[
\|\widehat{a}-\widehat{b}_{\alpha_{0}}\|_{\tau}< d_\tau \Big(\widehat{a},
\widehat{B_1(\overline{Q}^{S.O.T.})}\Big) + \epsilon.
\]
This implies that
\[
d_\tau\big(\widehat{a}, \widehat{B_1(Q)}\big)\leq
\|\widehat{a}-\widehat{b}_{\alpha_{0}}\|_{\tau} <d_\tau \Big(\widehat{a},
\widehat{B_1(\overline{Q}^{S.O.T.})}\Big)+ \epsilon.
\]
As $\epsilon>0$ was arbitrary, we get $d_\tau\Big(\widehat{a},
\widehat{B_1(Q)}\Big)\leq d_\tau\Big(\widehat{a},
\widehat{B_1(\overline{Q}^{S.O.T.})}\Big)$. This proves our assertion.

Then, we get
\[
\beta := \sup_{a\in B_1(P)}d_\tau \big(\widehat{a}, \widehat{B_1(Q)}\big)=
\sup_{a\in B_1(P)}d_\tau \Big(\widehat{a},
\widehat{B_1(\overline{Q}^{S.O.T.})}\Big)\]
\[
\leq \sup_{z\in
  B_1(\overline{P}^{S.O.T.})}d_\tau\Big(\widehat{z},
\widehat{B_1(\overline{Q}^{S.O.T.})}\Big)=: \alpha.
\] 
Further, for any  $\eta>0$, there exists a $ z_{0}\in
B_1\Big(\overline{P}^{S.O.T.}\Big)$ such that
\[
\alpha-\eta<
d_\tau \Big(\widehat{z_{0}}, \widehat{B_1(\overline{Q}^{S.O.T.})}\Big)\leq
\alpha.
\]
Then, by Kaplansky density theorem again, there exists a net
$\{z_{\lambda }\} \subset B_1(P)$ such that $z_{\lambda
}\stackrel{S.O.T.}{\longrightarrow} z_{0}$. In particular,
$\widehat{z_{\lambda}} {\longrightarrow} \widehat{z_{0}}$ with respect
to $\|\cdot\|_\tau$; and, since $d_\tau$ is continuous with respect to
$\|\cdot\|_\tau$, it follows that \( d_\tau \Big(\widehat{{z}_{\lambda }},
\widehat{B_1(\overline{Q}^{S.O.T.})}\Big)\rightarrow
d_\tau\Big(\widehat{z_{0}}, \widehat{B_1(\overline{Q}^{S.O.T.})}\Big).  \)
Thus, there exists a $\lambda_{0}$ such that \( d_\tau
\Big(\widehat{z}_{\lambda_{0}}, \widehat{B_1(\overline{Q}^{S.O.T.})}\Big)>\alpha-
\eta; \) so that, $\beta \geq \alpha-\eta$. Since $\eta > 0$ was
arbitrary, we have $\beta\geq \alpha$. Thus, $\alpha=\beta$, i.e., \(
d_{MT}(P, Q)= d_{MT}\Big({P}, \overline{Q}^{S.O.T.}\Big).  \) By a similar
argument, we conclude that
\[
d_{MT}(P, Q)= d_{MT}\Big(\overline{P}^{S.O.T.}, \overline{Q}^{S.O.T.}\Big).
\]
Rest follows from von Neumann's Double
Commutant theorem.
\end{proof}

\end{document}
